% Part: sets-functions-relations
% Chapter: sets
% Section: proofs-about-sets

\documentclass[../../../include/open-logic-section]{subfiles}

\begin{document}

\olfileid{sfr}{set}{prf}
\olsection{في إقامة البراهين على المجموعات}

\begin{editorial}
استُعيض عن هذا القسم بـ\verb|content/methods/proofs|، ولذلك لا يدخل في
هذا الفصل على الوضع الافتراضي.
\end{editorial}

\begin{explain}
تختصر لنا المجموعات والترميزات المتقدمة سبيل تعيين المجموعات وبيان
العلاقات بينها؛ غير أن الأمر كثيرًا ما يحتاج أيضًا إلى إثبات دعاوى في
تلك العلاقات. فإن لم تألف البرهان الرياضي، فنتدرج في مثال يسير خطوة
خطوة. نريد أن نثبت، لكل مجموعتين $X$ و$Y$، الهوية
$X \cap (X \cup Y) = X$. وسبيل إثبات هوية كهذه أن نرجع إلى أن المجموعة
لا تتعين إلا بـ!!{element}s المنتمية إليها: فمجموعتان متطابقتان إذا وفقط
إذا اشتركتا في !!{element}s نفسها. فيلزمنا هنا أولًا أن نثبت أن كل
!!{element} من $X \cap (X \cup Y)$ هو !!{element} من $X$، ثم أن نثبت
بالعكس أن كل !!{element} من $X$ هو !!{element} من
$X \cap (X \cup Y)$. أي نثبت (أ)
$X \cap (X \cup Y) \subseteq X$، ثم (ب) $X \subseteq X \cap (X \cup Y)$.

والدعوى العامة من نحو «كل !!{element}~$z$ من $X \cap (X
\cup Y)$ هو !!{element} من $X$» تثبت بأن نأخذ أولًا عنصرًا اعتباطيًا
$z \in X \cap (X \cup Y)$، ثم نستخرج من هذا الفرض~$z \in X$. ويسمى
هذا النمط أحيانًا «البرهان الشرطي العام». ونحتاج معه إلى استعمال
التعريفات الواردة في الفرض والنتيجة، وهما هنا تعريفَا «$\cap$» و«$\cup$».
فتقام حجة (أ) هكذا:

\begin{quote}
(أ) نثبت أولًا $X \cap (X \cup Y) \subseteq X$. ومعنى ذلك، بتعريف
$\subseteq$، أن $z \in X \cap (X \cup Y)$ يستلزم
$z \in X$، لكل~$z$. فافترض~$z \in X \cap (X \cup Y)$. ولما كان
$z$~!!{element} من تقاطع مجموعتين إذا وفقط إذا كان~!!{element} من
  % REVIEW-0721-N1: هذه خصوصية من اقتران ثابت، لا استدلال بطريق الأولى.
  كلتيهما، لزم~$z \in X$ ولزم أيضًا~$z \in X \cup Y$. وبخاصة
$z \in X$، وهو المطلوب.
\end{quote}

فتم الشطر الأول. وقد استعملنا في آخره أن لزوم الاقتران
($z \in X$ و$z \in X \cup Y$) عن فرض يوجب لزوم كل واحد من طرفيه عن
ذلك الفرض. وهذه هي قاعدة «حذف الاقتران»، أي حذف الرابط~$\land$. وننتقل
إلى (ب):

\begin{quote}
(ب) نثبت الآن $X \subseteq X \cap (X \cup Y)$؛ أي إن تعريف
$\subseteq$ يطلب أن يستلزم~$z \in X$ كون~$z \in X \cap
(X \cup Y)$، لكل~$z$. فافترض~$z \in X$. ولإثبات~$z \in X
\cap (X \cup Y)$، يكفينا بتعريف «$\cap$» إثبات (1)
$z \in X$، و(2) $z \in X \cup Y$. أما (1) فهو الفرض بعينه. وأما (2)،
فلأن~$z$ يكون~!!{element} من اتحاد مجموعات إذا وفقط إذا كان~!!{element}
من واحدة منها على الأقل، وكان~$z \in X$، وكان~$X \cup Y$ اتحاد~$X$
و$Y$، فقد تحقق المطلوب. ومن ثم~$z \in X \cup Y$. فثبت (1)
$z \in X$ و(2) $z \in X \cup Y$؛ ولذلك، بتعريف «$\cap$»،
$z \in X \cap (X \cup Y)$.
\end{quote}

وفي هذا تفصيل مقصود ليظهر وجه الاستدلال في المجموعات وعلاقاتها. أما في
العادة فنستغني عن كثير من التصريح، ولا نعيد كل تعريف. ولذلك يجيء البرهان
في متن أرفع مستوى موجزًا على الصورة الآتية.
\end{explain}


\begin{prop}[الامتصاص]
لكل مجموعتين $X$ و$Y$،
\[
X \cap (X \cup Y) = X
\]
\end{prop}

\begin{proof}
(أ) إذا كان $z \in X \cap (X \cup Y)$، لزم $z \in X$، ومن ثم
  $X \cap (X \cup Y) \subseteq X$.

(ب) وإذا كان $z \in X$، كان $z \in X \cup Y$ أيضًا، فكان
  $z \in X \cap (X \cup Y)$. ومن ثم
  $X \subseteq X \cap (X \cup Y)$.
\end{proof}

\begin{prob}
أثبت بالتفصيل أن $X \cup (X \cap Y) = X$، ثم صغ البرهان في عبارة
وجيزة. (تلميح: في اتجاه $X \cup (X \cap Y) \subseteq X$ تحتاج
إلى تقسيم البرهان إلى حالات، أي إلى قاعدة حذف الرابط~$\lor$.)
\end{prob}

\end{document}
